\section{Complete four-row multiplicities from cubic surfaces}

\subsection*{Finding A-038: the whole cubic four-row range and its padded comparison}
For every outer degree \(d\geq0\) and every partition
\(\mu\vdash3d\) with at most four parts,
\[
                    q^3_\mu(d[3])=K_3(\mu)=a_\mu(d[3]).
                                                               \tag{A-038.1}
\]
In particular, the entire proper partial-Hessian ray
\[
       \lambda_{3,4}=(8,3,2,2),\qquad D_4=5,
\]
has the exact actual multiplicities
\[
 q^3_{s\lambda_{3,4}}(5s[3])
       =K_3(s\lambda_{3,4})
       =a_{s\lambda_{3,4}}(5s[3])
                         \qquad(s\geq0).                 \tag{A-038.2}
\]
This is an equality with the complete ambient multiplicity,
not a claim that a chosen covariant algebra exhausts that
space, and not a closed formula for these plethysm coefficients.

There is also a comparison after arbitrary internal padding:
for every \(n\geq3\), \(d\geq0\), and
\(\mu\vdash nd\) with \(\ell(\mu)\leq4\),
\[
                         q^3_\mu(d[n])\leq K_n(\mu).
                                                               \tag{A-038.3}
\]
No second-row, third-row, or outer-degree bound is required.
Thus for every \(n\geq3\) and \(s\geq0\) the proper ray
\[
             \lambda_{n,4}=(5n-7,3,2,2),\qquad D_4=5,
\]
cannot be an actual multiplicity obstruction for the internally
padded \(3\)-by-\(3\) permanent, at any stretch.
The padded assertion (A-038.3) is an inequality; equality with
the full degree-\(n\) ambient space is not inferred.

\paragraph{The known geometric input.}
Beauville, \emph{Determinantal Hypersurfaces},
\emph{Michigan Math.\ J.}\ 48 (2000), 39--64,
Corollary~6.4, printed pp.~52--53, states that over an
algebraically closed field every smooth cubic surface has
a \(3\)-by-\(3\) linear determinantal equation. The immediately
preceding Proposition~6.2, p.~52, gives the associated
projectively normal curve and line-bundle resolution.
We use the existence part of Corollary~6.4, not its count of
equivalence classes. This is a classical theorem, not a
new determinantal representation result.

Fix a four-dimensional variable space \(U^*\). A smooth
cubic \(F\in\operatorname{Sym}^3U^*\) is therefore
\(\det M\), after absorbing a nonzero scalar into one row
of the linear matrix. Substitution of the entries of \(M\)
into the nine determinant variables is an endomorphism
substitution in the ambient space. Its image belongs to
the determinant orbit closure: the invertible substitutions
are dense in all endomorphisms, and the substitution map
is polynomial. Smooth cubics form a nonempty open dense
subset, as the Fermat cubic verifies. Closedness now gives
\[
          \operatorname{Sym}^3U^*
                    \subseteq\overline{\mathrm{GL}_9\det_3}.
                                                               \tag{A-038.4}
\]

\paragraph{One zero entry converts these determinants to permanents.}
Let \(M(x)=\sum_{i=1}^4x_iA_i\) be any \(3\)-by-\(3\)
matrix of four-variable linear forms. There exist nonzero
column vectors \(u,v\in\mathbb C^3\) such that
\[
                         u^{\mathsf T}A_i v=0
                                      \quad(1\leq i\leq4).
                                                               \tag{A-038.5}
\]
Indeed the affine cone of rank-at-most-one \(3\)-by-\(3\)
matrices has dimension \(3+3-1=5\). Impose the four
homogeneous linear equations obtained by pairing a
rank-one matrix with the \(A_i\). The height bound for an
ideal generated by four elements gives an intersection of
dimension at least one. Thus it contains a nonzero rank-one
matrix, which factors as \(uv^{\mathsf T}\), with the
transpose convention chosen to give (A-038.5).
Equivalently, four hyperplanes cannot miss the
four-dimensional Segre variety in \(\mathbb P^8\).

Complete \(u^{\mathsf T}\) to the first row of an invertible
matrix \(P\), and \(v\) to the first column of an invertible
matrix \(Q\). Then \((PMQ)_{11}=0\). For arbitrary entries
the exact identity
\[
 \det\begin{pmatrix}0&a&b\\c&d&e\\f&g&h\end{pmatrix}
     =
 \operatorname{per}
       \begin{pmatrix}0&a&b\\c&-d&e\\f&g&-h\end{pmatrix}
     =-ach+aef+bcg-bdf                                      \tag{A-038.6}
\]
holds. Rescaling one row removes the scalar
\(\det(P)\det(Q)\).
Consequently every smooth quaternary cubic is a linear
substitution into \(\operatorname{per}_3\) as well.
The same endomorphism and closedness argument proves
\[
          \operatorname{Sym}^3U^*
                    \subseteq\overline{\mathrm{GL}_9\operatorname{per}_3}.
                                                               \tag{A-038.7}
\]
The dimension argument is used only for the zero entry;
it is not being used to infer dominance of a coefficient
map from a parameter count.

\begin{proof}[Complete multiplicities at inner degree three]
Apply the subspace criterion O-006.a to each of
(A-038.4) and (A-038.7). To recall its all-copy content,
the degree-\(d\) polynomial coordinate space has the
functorial decomposition
\[
 \operatorname{Sym}^d\operatorname{Sym}^3V
      =\bigoplus_{\mu\vdash3d}S_\mu V\otimes M_\mu.
\]
An equivariant ideal is
\(\bigoplus_\mu S_\mu V\otimes J_\mu\), where
\(J_\mu\subseteq M_\mu\).
Restriction to a four-variable subspace acts as
\(S_\mu(V\to U)\otimes\mathrm{id}_{M_\mu}\).
For \(\ell(\mu)\leq4\) the Schur map is surjective and
its target is nonzero. If the variety contains the full
four-variable form space, no nonzero \(J_\mu\) can
restrict to zero there. Thus \(J_\mu=0\).
Both varieties retain all \(a_\mu(d[3])\) copies, proving
(A-038.1), and substitution of the stated ray proves
(A-038.2).
\end{proof}

\paragraph{The entire padded cubic family is available on four variables.}
Let \(n\geq3\). For every linear form \(L\in U^*\) and
every cubic \(F\in\operatorname{Sym}^3U^*\),
\[
                      L^{\,n-3}F\in
                            \overline{\mathrm{GL}_{n^2}\det_n}.
                                                               \tag{A-038.8}
\]
For smooth \(F=\det M\) this is the explicit identity
\[
        L^{\,n-3}F
           =\det\begin{pmatrix}
                        M&0\\0&L I_{n-3}
                 \end{pmatrix}.                              \tag{A-038.9}
\]
For fixed \(L\) pass to the closure of the smooth cubics
to obtain every \(F\); this also includes \(L=0\).
When \(n=3\) the lower block is empty and the factor
\(L^0\) is one.

\begin{proof}[The padded comparison includes every highest-weight copy]
Use the internal convention
\[
              P_n=X_{11}^{\,n-3}\operatorname{per}_3,
\]
with the nine entries among the \(n^2\) ambient variables.
Restrict any endomorphism substitution of \(P_n\) to the
first four variables. It has the form
\[
                          L^{\,n-3}F,
\]
where \(L\) is the restricted substituted \(X_{11}\) and
\(F\) is a four-variable cubic. The dependence of \(F\)
on that same distinguished entry causes no problem:
(A-038.8) holds for \emph{every} pair \((L,F)\).
Since the four-variable determinant locus is closed, the
restriction of every point of the permanent orbit closure
also belongs to it.

Let \(H_\mu\) be the complete ambient highest-weight
multiplicity space of a type with \(\ell(\mu)\leq4\).
Every polynomial in \(H_\mu\) involves only coefficient
generators supported on the first four variables: its
torus weights in all later coordinates are zero, and all
generator weights are nonnegative.
If \(w\in H_\mu\) vanishes on the determinant closure,
its value on a permanent point is therefore its value
on that point's four-variable restriction. This restriction
belongs to the determinant locus by (A-038.8), so the
value is zero. We have proved the inclusion of complete
highest-weight ideal kernels
\[
       I(\mathrm{Det}_n)_d\cap H_\mu
                    \subseteq I(\mathrm{Per}_{n,3})_d\cap H_\mu.
                                                               \tag{A-038.10}
\]
Subtract their dimensions from \(\dim H_\mu=a_\mu(d[n])\).
This gives \(q^3_\mu(d[n])\leq K_n(\mu)\), proving
(A-038.3). No equality or isomorphism between the two
full coordinate rings is asserted.
\end{proof}

\paragraph{Scope and geometric dictionary.}
The cubic-surface theorem gives a whole linear subspace in
each cubic orbit closure. The block construction gives the
image of the entire multiplication morphism
\(U^*\times\operatorname{Sym}^3U^*\to\operatorname{Sym}^nU^*\),
\((L,F)\mapsto L^{n-3}F\), in the determinant closure.
The conclusion is an inclusion of the actual ideal kernels,
which controls all multiplicity copies, including vectors
outside the binary form/partial-Hessian algebra.
It does not extend the cubic-surface theorem to general
quartic surfaces, and does not imply \(K_n=a_\mu\) for
all four-row types when \(n\geq4\); O-008 disproves that
last assertion.

The torus-weight and coherent representation interpretation
is SGA~3, I, Propositions~4.7.2--4.7.3, pp.~34--36.
The Cauchy and Schur operations are those of
SGA~6, VI, Theorem~3.3, pp.~392--393, and Macdonald,
I(4.3), p.~63, and Appendix~A(5.3)--(5.6), pp.~155--156.
The rank-one cone is the cone over the Segre image of
\(\mathbb P^2\times\mathbb P^2\); the four-equation
height bound is applied to its coordinate ring over
\(\mathbb C\).
The known representation theorem was read in the author's
published PDF at
\url{https://math.univ-cotedazur.fr/u/beauvill/pubs/det.pdf},
with the exact original-page reading and file hash
recorded in the source check log. The source PDF and
rendered evidence remain outside the work package.
